\chapter{Il tessuto connettivo specializzato}

% ---------------------------------------------------------------------------
\section{Il tessuto cartilagineo}

\subsection{Generalità}

La \textbf{cartilagine} è un materiale filogeneticamente molto antico: nei
pesci cartilaginei (ad esempio gli squali) costituisce l'intero scheletro,
mentre nell'uomo persiste in tutte quelle sedi in cui è necessaria una
struttura solida ma al tempo stesso deformabile e resistente alla
compressione. Si ritrova infatti nel setto nasale, nella laringe e
nell'apparato respiratorio, nei dischi intervertebrali, nel padiglione
auricolare, nelle articolazioni di spalla, gomito e testa omerale, nei
menischi, nella sinfisi pubica e nelle articolazioni di mano e piede.

Come il tessuto osseo, la cartilagine è priva di vasi sanguigni: il suo
nutrimento avviene per diffusione a partire dal \textbf{pericondrio}, lo
strato di connettivo denso che la riveste. Questa assenza di
vascolarizzazione rende la \textbf{rigenerazione} della cartilagine un
processo molto lento e spesso incompleto.

\subsection{I condrociti}

Le cellule della cartilagine, i \textbf{condrociti}, hanno una caratteristica
forma \textbf{globosa} e sono alloggiate in piccole cavità della matrice dette
\textit{lacune}. Essendo il tessuto cartilagineo non
irrorato, i condrociti traggono nutrimento esclusivamente per diffusione
attraverso la matrice extracellulare, il che limita fortemente la loro
capacità proliferativa e, di conseguenza, la capacità dell'intero tessuto di
rigenerarsi dopo un danno.

Si distinguono tre tipi di cartilagine, in base alla composizione della
matrice extracellulare: \textbf{ialina}, \textbf{fibrosa} ed
\textbf{elastica}.

\subsection{Cartilagine ialina}

La \textbf{cartilagine ialina} è il tipo più diffuso: costituisce le
superfici articolari delle ossa lunghe, le cartilagini costali, la maggior
parte dello scheletro laringeo e gli anelli di trachea e bronchi. È inoltre
il tessuto di cui è costituito il \textbf{modello cartilagineo} da cui, nella
vita prenatale, si sviluppa la maggior parte dello scheletro per
\textbf{ossificazione endocondrale}. La sua matrice, otticamente omogenea al
microscopio ottico, è ricca di collagene di tipo II e di proteoglicani, che le
conferiscono resistenza alla compressione e un basso coefficiente di attrito
sulle superfici articolari.

\subsection{Cartilagine fibrosa}

La \textbf{cartilagine fibrosa} contiene componenti fibrose visibili al
microscopio ottico ed è una forma di transizione tra il tessuto connettivo
denso e la cartilagine propriamente detta: è caratterizzata da grossi fasci
fibrosi immersi in una scarsa matrice cartilaginea, contenente quantità
variabili di proteoglicani. Si trova nei dischi
interarticolari, in particolare nei \textbf{menischi}, dove risulta
generalmente ben fissata all'osso circostante ed è più resistente alla
trazione rispetto alla cartilagine ialina.

\subsection{Cartilagine elastica}

La \textbf{cartilagine elastica} è molto rara nell'organismo. È presente nel
padiglione auricolare e nella tuba uditiva (di Eustachio), ai quali conferisce
\textbf{deformabilità}: è infatti capace di resistere a urti improvvisi senza
rompersi e di flettersi grazie alla presenza di numerose fibre elastiche
nella sua matrice extracellulare.

% ---------------------------------------------------------------------------
\section{Il tessuto osseo}

\subsection{Generalità}

L'\textbf{osso} è una parte anatomica solida dei vertebrati; l'insieme delle
ossa costituisce lo \textbf{scheletro}. La straordinaria comprimibilità e
resistenza alle forze di taglio dell'osso, unite al suo peso relativamente
modesto, lo rendono il materiale ideale per lo scheletro dei grandi animali
terrestri, incluso l'uomo.

Alla nascita lo scheletro umano presenta circa \textbf{270 ossa}; nell'adulto
si riducono a \textbf{206}, legate tra loro da 68 articolazioni, poiché
durante lo sviluppo alcune ossa si fondono tra loro formandone una sola. Il
peso dello scheletro corrisponde a circa il \textbf{20\,\%} del peso
corporeo.

Il tessuto osseo è un \textbf{connettivo specializzato}, la cui matrice
extracellulare è calcificata e imprigiona le cellule che l'hanno prodotta.
Nonostante la sua durezza, è un tessuto dinamico, in grado di rimodellarsi
continuamente in relazione alle forze che agiscono su di esso. Ogni osso è
rivestito esternamente dal \textbf{periostio} e delimitato internamente
dall'\textbf{endostio}; la cavità centrale contiene il midollo osseo, sede
dell'emopoiesi.

\subsection{Funzioni del tessuto osseo}

Le funzioni fisiologiche delle ossa si possono raggruppare in tre categorie:
\textbf{meccaniche}, \textbf{emopoietiche} e \textbf{metaboliche}.

\subsubsection{Funzioni meccaniche}

\begin{itemize}
  \item \textbf{Protettiva}: protegge gli organi interni (il cranio per
    l'encefalo, la gabbia toracica per cuore e polmoni).
  \item \textbf{Motoria}: fornisce le leve per il movimento nello spazio
    tridimensionale.
  \item \textbf{Sensoriale}: gli ossicini dell'udito trasmettono il suono agli
    organi interni dell'orecchio.
  \item \textbf{Sostenitiva}: costituisce l'impalcatura del corpo, fornendo
    attacco e sostegno a tutti i tessuti molli.
\end{itemize}

\subsubsection{Funzioni emopoietiche}

Il midollo delle ossa lunghe e il tessuto spugnoso delle ossa piatte
contengono le cellule staminali emopoietiche, che generano gli eritrociti e i
leucociti. Il midollo osseo si distingue in tre tipi:

\begin{itemize}
  \item \textbf{Midollo rosso}: presente nelle ossa del neonato e, nell'adulto,
    nel tessuto spugnoso delle ossa brevi e piatte e nell'epifisi vicina a
    femore e omero; è la sede attiva dell'emopoiesi.
  \item \textbf{Midollo giallo}: presente nella diafisi delle ossa lunghe
    dell'adulto, è ricco di cellule adipose e rappresenta una riserva
    energetica.
  \item \textbf{Midollo grigio}: presente nell'anziano, è ormai inattivo dal
    punto di vista emopoietico.
\end{itemize}

\subsubsection{Funzioni metaboliche}

\begin{itemize}
  \item \textbf{Riserva di minerali}: calcio e fosforo.
  \item \textbf{Riserva di grassi}: il midollo giallo contiene acidi grassi
    prelevabili dal sangue in caso di necessità.
  \item \textbf{Riserva di fattori di crescita}: la matrice ossea mineralizzata
    ne contiene quantità importanti (ad esempio la proteina morfogenetica
    ossea), rilasciati localmente in caso di frattura per innescare e
    accelerare la guarigione.
  \item \textbf{Detossificazione}: la componente inorganica dell'osso può
    assorbire metalli pesanti e altri elementi estranei, rimuovendoli dal
    circolo sanguigno.
  \item \textbf{Equilibrio acido-base}: l'osso tampona le variazioni di pH del
    sangue assorbendo o rilasciando sali minerali e ioni.
  \item \textbf{Secrezione endocrina}: l'osso contribuisce al controllo del
    metabolismo del fosforo.
\end{itemize}

\subsection{Osso compatto e osso spugnoso}

Il tessuto osseo si presenta in due forme: l'\textbf{osso compatto}, che
costituisce la porzione esterna delle ossa lunghe, e l'\textbf{osso
spugnoso}, che ne occupa le epifisi e la porzione interna delle ossa piatte
e brevi.

L'\textbf{osso compatto} è un tessuto solido, privo di spazi se non quelli
occupati dalle cellule, dai loro processi e dai vasi sanguigni. È formato da
\textbf{colonne ossee} (osteoni) parallele all'asse maggiore dell'osso,
ciascuna delle quali è costituita da strati concentrici, o \textbf{lamelle}
(\textit{osso lamellare}), disposti attorno a un canale centrale.

La \textbf{morfologia dell'osso compatto} è definita dai seguenti elementi:

\begin{itemize}
  \item \textbf{Canale di Havers}: canale centrale che accoglie vasi
    sanguigni e nervi.
  \item \textbf{Canale di Volkmann}: canale trasversale che collega tra loro
    i canali di Havers di osteoni adiacenti.
  \item \textbf{Lamelle}: cerchi concentrici originati dalla deposizione
    apposizionale dell'osso.
  \item \textbf{Lacuna}: piccola cavità che ospita l'osteocita.
  \item \textbf{Canalicoli}: sottili canali che accolgono i processi degli
    osteociti, permettendo lo scambio di metaboliti tra cellule adiacenti.
\end{itemize}

L'\textbf{osso spugnoso} è invece composto da una rete di \textbf{trabecole}
ossee sottili, separate da un labirinto di spazi che accolgono il midollo
osseo. Le trabecole dipartono dalle lamelle
circonferenziali interne verso la cavità midollare e sono costituite da
lamelle ossee irregolari con lacune contenenti osteociti. A differenza
dell'osso compatto, l'osso spugnoso non possiede sistemi di Havers: gli
osteociti scambiano i metaboliti con i sinusoidi del midollo attraverso i
canalicoli. Le trabecole sono rivestite dall'\textbf{endostio}, un sottile
strato di connettivo che contiene cellule osteoprogenitrici, osteoblasti e
osteoclasti.

Al microscopio elettronico a scansione la differenza tra le due forme di
tessuto osseo risulta evidente: la superficie compatta e densa dell'osso
compatto, percorsa dai canali di Havers, si contrappone alla struttura
reticolare e porosa dell'osso spugnoso.

\subsection{Periostio ed endostio}

Il \textbf{periostio} è un sottile strato di connettivo denso irregolare,
ricco di collagene, che circonda l'osso. Si inserisce sull'osso tramite le \textbf{fibre di
Sharpey}, fasci di fibre collagene che penetrano nella matrice ossea. Si
distinguono due strati:

\begin{itemize}
  \item uno \textbf{strato fibroso esterno}, che veicola vasi e nervi;
  \item uno \textbf{strato cellulare interno}, che contiene le cellule
    osteoprogenitrici, capaci di differenziarsi in osteoblasti.
\end{itemize}

L'\textbf{endostio} tappezza tutte le superfici interne delle cavità ossee
(cavità midollare, trabecole dell'osso spugnoso, canali di Havers e di
Volkmann). È composto da un unico strato di cellule osteoprogenitrici, più
sottile del periostio, ed è deputato al nutrimento e alla produzione di
nuove cellule ossee.

\subsection{Ossa piatte}

Le \textbf{ossa piatte} della volta cranica si sviluppano con modalità
diverse dalle ossa lunghe: sono costituite da un \textbf{tavolato interno} e
un \textbf{tavolato esterno} di osso compatto, che racchiudono uno strato
centrale di osso spugnoso detto \textbf{diploe}. Il tavolato esterno è
rivestito dal \textbf{pericranio}, mentre il tavolato interno è rivestito
dalla \textbf{dura madre}, che riveste anche il periostio e protegge
l'encefalo.

\subsection{Osteogenesi}

Durante la vita prenatale l'osteogenesi (ossificazione) si sviluppa
generalmente in modo indiretto, passando attraverso un \textbf{modello
cartilagineo}: si parla in questo caso di \textbf{ossificazione
endocondrale}. Fanno eccezione le ossa piatte della volta cranica, che
ossificano direttamente a partire da un abbozzo connettivale, senza passare
attraverso una fase cartilaginea (\textbf{ossificazione
intramembranosa}).

\subsubsection{Ossificazione endocondrale}

Nell'embrione, nella sede in cui si formerà l'osso, è presente un abbozzo
cartilagineo che si accresce sia per crescita interstiziale sia per crescita
apposizionale. A partire dalla 6\textsuperscript{a}--9\textsuperscript{a}
settimana di vita fetale circa, la cartilagine viene gradualmente demolita e
sostituita da un manicotto osseo; il processo ha inizio nelle ossa lunghe a
partire dal centro della diafisi (\textbf{centro primario di
ossificazione}).

I condrociti della regione centrale si ipertrofizzano, accumulano glicogeno
nel citoplasma e si vacuolizzano: l'ipertrofia porta a un aumento del volume
della lacuna, con riduzione della matrice circostante, che inizia a
calcificare. Alla giunzione tra cartilagine e osso si formano così
\textbf{complessi calcificati cartilagine/osso} in via di riassorbimento,
ben visibili al microscopio ottico.

Successivamente compaiono, con modalità simili, i \textbf{centri secondari
di ossificazione} nelle epifisi: a differenza del centro primario, qui non si
forma il manicotto osseo. Le cellule osteoprogenitrici invadono la
cartilagine epifisaria, si differenziano in osteoblasti e cominciano a
secernere la matrice ossea. L'ossificazione procede parallelamente a quella
diafisaria, ma non coinvolge le superfici articolari né la piastra
epifisaria (cartilagine di accrescimento).

Al momento della nascita, nelle zone epifisarie delle ossa lunghe compaiono
i nuclei di ossificazione, che si estendono perifericamente; la cartilagine
originaria rimane limitata alle superfici articolari e a una zona compresa
tra epifisi e diafisi. Per effetto degli ormoni sessuali, durante la
pubertà le commissure epifisarie si ossificano, portando a termine la
crescita in lunghezza dell'osso: la massa ossea raggiunge così il proprio
picco, per poi ridursi progressivamente con l'età, in modo più marcato nella
donna dopo la menopausa.

Lo sviluppo di un osso lungo procede quindi dall'embrione al bambino con la
progressiva sostituzione del modello cartilagineo da parte del tessuto osseo,
a partire dal centro diafisario.

\subsubsection{Osteogenesi e vitamina D}

La \textbf{vitamina D} è di importanza fondamentale per l'osteogenesi,
poiché influenza il bilancio del calcio. Nei pazienti lungodegenti, una
insufficiente esposizione alla radiazione UV necessaria alla sua sintesi
cutanea può determinarne una carenza; nella profilassi dell'osteoporosi la
vitamina D può essere somministrata farmacologicamente, mentre
nell'alimentazione infantile è necessario garantirne un apporto sufficiente,
insieme al calcio, per favorire l'osteogenesi. Pochi alimenti ne contengono
quantità apprezzabili: tra questi l'olio di fegato di merluzzo, i pesci
grassi (salmone, aringa), il latte e i suoi derivati, le uova, il fegato e le
verdure a foglia verde.

\subsubsection{Osteogenesi imperfetta}

L'\textbf{osteogenesi imperfetta} è una malattia genetica rara, con
distribuzione uguale tra i sessi (M/F = 1:1), che causa gravi problemi a
carico dello scheletro, clinicamente manifesti come \textbf{fragilità
ossea}. Ne sono note 7 varianti a diversa gravità, raggruppate in 4 tipi
principali:

\begin{itemize}
  \item \textbf{Tipo I}: ritardo dell'accrescimento, fratture ossee, cifosi,
    scoliosi, perdita dell'udito.
  \item \textbf{Tipo II}: fatale durante la vita intrauterina, con fratture
    multiple che si manifestano quando il feto è ancora in utero.
  \item \textbf{Tipo III}: fratture alla nascita, deformazioni progressive,
    bassa statura.
  \item \textbf{Tipo IV}: clinicamente meno grave, statura poco ridotta,
    fragilità ossea lieve o moderata, deformità variabili.
\end{itemize}

\subsection{Le cellule dell'osso}

Il tessuto osseo comprende quattro tipi cellulari principali: le cellule
\textbf{osteoprogenitrici}, gli \textbf{osteoblasti}, gli \textbf{osteociti}
e gli \textbf{osteoclasti}.

\subsubsection{Cellule osteoprogenitrici}

Derivano dalle cellule mesenchimali embrionali e mantengono la capacità di
dividersi. Si trovano nello strato interno del periostio e nell'endostio, da
cui, per differenziamento, originano gli osteoblasti.

\subsubsection{Osteoblasti}

Gli \textbf{osteoblasti} sintetizzano e secernono i costituenti organici
della matrice ossea. Sono cellule di forma cubica o cilindrica, disposte in
superficie, con nucleo eccentrico opposto ai granuli di secrezione; il
reticolo endoplasmatico rugoso e l'apparato di Golgi sono molto sviluppati,
coerentemente con l'intensa attività secretoria. Processi citoplasmatici li
mettono in contatto con le cellule vicine tramite giunzioni comunicanti.

Gli osteoblasti che, in superficie, cessano di produrre matrice ritornano a
uno stato quiescente e vengono chiamati \textbf{cellule delimitanti l'osso}:
sono simili alle cellule osteoprogenitrici ma incapaci di dividersi, pur
potendo riprendere l'attività secretoria.

\subsubsection{Osteociti}

Gli \textbf{osteociti} sono le cellule mature dell'osso, presenti in numero
di 20\,000--30\,000 per mm\textsuperscript{3}, derivate dagli osteoblasti che
restano intrappolati nelle lacune della matrice che essi stessi hanno
prodotto. Sono più piccoli e meno basofili degli osteoblasti, con forma
adattata a quella della lacuna, nucleo appiattito, pochi organuli e scarso
reticolo endoplasmatico rugoso. Possiedono processi citoplasmatici
inter-connettenti che decorrono nei canalicoli, uniti da giunzioni
comunicanti (gap-junction).

Pur apparendo cellule quiescenti, gli osteociti secernono sostanze che
mantengono la matrice ossea in buono stato, rilasciano ioni calcio dalla
matrice quando la richiesta corporea aumenta e facilitano l'attività dei
preosteoblasti nel rimodellamento dello scheletro.

\subsubsection{Osteoclasti}

Gli \textbf{osteoclasti} sono cellule multinucleate che derivano da
progenitori macrofagico-monocitari del midollo osseo, in seguito alla
stimolazione di fattori specifici. Si muovono sulla
superficie dell'osso e ne riassorbono la matrice sotto lo stimolo
dell'\textbf{ormone paratiroideo} e della \textbf{calcitonina}, occupandosi
della decalcificazione, della digestione e dell'assorbimento dei materiali
digeriti. Occupano piccole depressioni scavate nella superficie ossea in via
di riassorbimento, dette \textbf{lacune di Howship}.

\subsection{L'equilibrio tra osteoblasti e osteoclasti}

L'equilibrio tra riassorbimento e neoformazione è vitale per l'osso:
l'osso distrutto dagli osteoclasti deve essere costantemente rimpiazzato da
nuovo osso prodotto dagli osteoblasti. Tutto il processo è regolato da
cellule, e può alterarsi quando le cellule invecchiano o i sistemi di
controllo si sbilanciano. L'osso è inoltre sempre alla mercé delle necessità
corporee di calcio, che hanno priorità per il funzionamento di nervo e
muscolo: per questo motivo il carico meccanico quotidiano è essenziale al
mantenimento della sua massa e della sua struttura.
